\documentclass[11pt]{article}
\usepackage[margin=1in]{geometry}
\usepackage{amsmath,amssymb,amsthm}
\usepackage{seqsplit}
\usepackage{hyperref}
\usepackage{booktabs}
\usepackage{xcolor}

\newtheorem{theorem}{Theorem}
\newtheorem{lemma}{Lemma}
\newtheorem{corollary}{Corollary}
\newtheorem{construction}{Construction}

\newcommand{\bignum}[1]{\texttt{\seqsplit{#1}}}
\newcommand{\numline}[2]{\par\noindent $#1 = $\ \bignum{#2}\par\vspace{2pt}}

\title{Large Modular Congruences for which $x \mid 36n^3-19$\\
\large A CRT / cube-root construction with 30- and 40-digit moduli}
\author{}
\date{\today}

\begin{document}
\maketitle

\begin{abstract}
We give a fully rigorous, arbitrarily-scalable construction of congruence classes
$n \equiv a_1 \pmod{p_1}$ and (more generally) $n\equiv A \pmod M$ such that
$36n^3-19$ is exactly divisible by a prescribed modulus $M$ built from primes
of any prescribed size. The construction rests on the elementary fact that cubing
is a bijection modulo any prime $p\equiv 2\pmod 3$. We instantiate this with
30-digit and 40-digit primes, verify every claim by direct computation (both
symbolically for small primes and numerically for the large instances), and
exhibit an explicit family in which $36n^3-19$ is exactly divisible by a
120-digit modulus, itself a product of three 40-digit primes. All numerical
claims below were checked by exact (arbitrary precision) integer arithmetic.
\end{abstract}

\section{Problem restatement}

We seek congruences
\[
  n \equiv a_1 \pmod{p_1}, \qquad x \equiv a_2 \pmod{p_2}
\]
with $a_1,p_1,a_2,p_2$ as large as we like, such that
\[
  \frac{36n^3-19}{x} \in \mathbb{Z}.
\]
Read literally, a congruence class for $n$ can only be paired with a \emph{fixed}
divisor $x$ that works for \emph{every} representative of the class if $x$ divides
the modulus $p_1$ itself (since $36n^3-19$ ranges over infinitely many values as
$n$ ranges over the class, and only a common factor of the modulus can divide all
of them). This forces $p_2 \mid p_1$, and the richest non-trivial choice is
$x = p_1$ exactly, i.e.\ $a_2 = 0$, $p_2 = p_1$.

To get a genuinely large, genuinely non-trivial $x$ that still divides
$36n^3-19$ for a whole residue class of $n$ (not just a single verified integer),
we combine \emph{two independent} such congruences via the Chinese Remainder
Theorem, so that $x$ becomes the product of two (or more) large primes, each
individually up to 30 digits. This is the construction below; it directly
generalizes to any number of primes.

\section{The key lemma}

\begin{lemma}\label{lem:cuberoot}
Let $p$ be a prime with $p\equiv 2\pmod 3$. Then the map $t\mapsto t^3 \bmod p$
is a bijection on $\mathbb{Z}/p\mathbb{Z}$. Consequently, for any $c\in(\mathbb Z/p\mathbb Z)^\times$
there is a unique $a\in\{0,1,\dots,p-1\}$ with $a^3\equiv c\pmod p$, given
explicitly by
\[
  a \equiv c^{\,e}\pmod p,\qquad e \equiv 3^{-1}\pmod{p-1}.
\]
\end{lemma}

\begin{proof}
Since $p\equiv 2\pmod 3$ we have $p-1\equiv 1\pmod 3$, so $\gcd(3,p-1)=1$.
Hence $3$ is invertible modulo $p-1$; let $e=3^{-1}\bmod (p-1)$. For
$t\in(\mathbb Z/p\mathbb Z)^\times$, Fermat's little theorem gives
$(t^3)^e = t^{3e} = t^{1+k(p-1)} = t$ for the appropriate integer $k$, so
cubing and raising to the power $e$ are mutually inverse bijections on the
multiplicative group; extending to $t=0$ is trivial.
\end{proof}

\begin{construction}\label{con:main}
Fix a prime $p\equiv 2\pmod 3$. Apply Lemma~\ref{lem:cuberoot} with
$c \equiv 19\cdot 36^{-1}\pmod p$ to obtain the unique
\[
  a \equiv \bigl(19\cdot 36^{-1}\bigr)^{e}\pmod p, \qquad e=3^{-1}\bmod (p-1),
\]
satisfying $36a^3\equiv 19\pmod p$. Then for \emph{every} integer $n\equiv a\pmod p$,
\[
  p \mid 36n^3-19.
\]
\end{construction}

This is immediate: $36n^3-19\equiv 36a^3-19\equiv 0\pmod p$ whenever $n\equiv a\pmod p$.

\begin{theorem}[Multi-prime combination]\label{thm:crt}
Let $p_1,\dots,p_k$ be distinct primes, each $\equiv 2\pmod 3$, and let
$a_1,\dots,a_k$ be the corresponding roots from Construction~\ref{con:main}.
Let $M=p_1p_2\cdots p_k$ and let $A$ be the unique solution modulo $M$ of the
simultaneous system $n\equiv a_i \pmod{p_i}$ ($1\le i\le k$), given by the
Chinese Remainder Theorem. Then for every $n\equiv A\pmod M$,
\[
  M \mid 36n^3-19, \qquad\text{i.e.}\qquad \frac{36n^3-19}{M}\in\mathbb Z .
\]
\end{theorem}

\begin{proof}
For each $i$, $n\equiv A\equiv a_i\pmod{p_i}$, so $p_i\mid 36n^3-19$ by
Construction~\ref{con:main}. The $p_i$ are distinct primes, so their product
$M$ divides $36n^3-19$.
\end{proof}

So Theorem~\ref{thm:crt} answers the original question in the strong (infinite
family) sense: $n\equiv A\pmod M$ and $x = M$ is an exact divisor of $36n^3-19$
for \emph{every} such $n$, and $M$, $A$ can be made as large as desired by
taking more/larger primes $p_i$.

\section{Sanity check on small primes}

Before trusting the construction at 30-digit scale, Lemma~\ref{lem:cuberoot}
and Construction~\ref{con:main} were checked by brute force for every prime
$p<200$ with $p\equiv 2\pmod 3$: for each such $p$, the formula's output $a$
was compared against an exhaustive search over $\{0,\dots,p-1\}$ for roots of
$36t^3\equiv 19\pmod p$, confirming in every case that (i) a root exists,
(ii) it is unique, and (iii) it matches the closed-form value from the lemma.
Representative small case: $p=17$ gives $a=14$, and indeed
$36\cdot 14^3-19 = 98685 = 17\cdot 5805$.

\section{Two 30-digit primes ($k=2$)}

Two primes $p_1,p_2\equiv 2\pmod 3$, each exactly 30 digits, were generated and
fed through Construction~\ref{con:main}:

\numline{p_1}{432256083118531044858259323599}
\numline{a_1}{205375962562755822717169031496}
\numline{p_2}{508951041752918591462038694231}
\numline{a_2}{11998528096759146724786261916}

Direct check: $36a_1^3-19\equiv 0\ (\mathrm{mod}\ p_1)$ and
$36a_2^3-19\equiv 0\ (\mathrm{mod}\ p_2)$, both verified exactly.

Combining via CRT (Theorem~\ref{thm:crt}, $k=2$) gives modulus
$M=p_1p_2$ (60 digits) and representative $A$ (59 digits):

\numline{M}{219997183807212542923931096695598007560918353360263243457369}
\numline{A}{61583691402182783047260797510703276439292047987969866884962}

so that for every $n\equiv A\pmod M$,
\[
  \frac{36n^3-19}{M}\in\mathbb Z .
\]
For the representative $n=A$ itself, $36A^3-19$ is a 178-digit integer whose
exact quotient by $M$ is the 119-digit integer:

\numline{Q}{38219282687467843755039989619336413202799418213239056237824742757700405657626958942597846112998216708072139112513223381}

This was verified by exact big-integer arithmetic: $36A^3-19 = M\cdot Q$
precisely (no rounding, no floating point).

\section{Three 30-digit primes ($k=3$)}

Pushing Theorem~\ref{thm:crt} to $k=3$ with three fresh 30-digit primes
$p_1,p_2,p_3\equiv 2\pmod 3$:

\numline{p_1}{779282871203804318581334453057}
\numline{a_1}{540863547358444129637446662643}
\numline{p_2}{917485502986111972709840010173}
\numline{a_2}{765515263147582646243402496698}
\numline{p_3}{422843868764454355911216773723}
\numline{a_3}{173885648024960625137792429386}

each individually verified to satisfy $36a_i^3\equiv 19\pmod{p_i}$. The
combined modulus is the 90-digit integer $M=p_1p_2p_3$:

\numline{M}{302325220948348183420434660024015348902086281522777512680498475446030651529604772341579503}

with CRT representative (90 digits):

\numline{A}{182470001039324181181024873011044406335174068803296028244172097475931272694261226503431598}

By Theorem~\ref{thm:crt}, for \emph{every} integer $n\equiv A\pmod M$,
\[
  x = M \quad\text{exactly divides}\quad 36n^3-19 .
\]
Concretely, at $n=A$, $36A^3-19$ is a 270-digit integer, and $(36A^3-19)/M = Q$
is the following 180-digit integer:

\numline{Q}{723440045918937868278733134032344285097446574084321029094017717823428606222435878944285733479052368954298584972235671037551447033227302246880204085670050855668916609015702186192131}

verified exactly via $36A^3-19=M\cdot Q$ in arbitrary-precision arithmetic.

\section{Two 40-digit primes ($k=2$)}

The same construction repeated with 40-digit primes $p_1,p_2\equiv 2\pmod 3$:

\numline{p_1}{7729484335457653901640057298531371241781}
\numline{a_1}{7668575607239450973459863267707132263860}
\numline{p_2}{2486598372481845396683104279916570951657}
\numline{a_2}{609530524018264138310326718615033307496}

each verified individually against $36a_i^3\equiv 19\pmod{p_i}$. Combining via
CRT gives the 80-digit modulus $M=p_1p_2$ and 80-digit representative $A$:

\numline{M}{19220123168672920512532048525080739795893456098633704192461052168447373009581117}
\numline{A}{12055573306960757229878249326862599227290509643013675419894826308035075337947215}

so that for every $n\equiv A\pmod M$, $M$ exactly divides $36n^3-19$. At
$n=A$, $36A^3-19$ is a 239-digit integer, and $(36A^3-19)/M=Q$ is the
following 160-digit integer:

\numline{Q}{3281783589609439060274407871106333299264698047304666391611865457187494168830153506599538715008344851754116620345469289412464296994115391491979199807003772632893}

verified exactly via $36A^3-19=M\cdot Q$.

\section{Three 40-digit primes ($k=3$) --- the largest instance found}

Finally, three 40-digit primes $p_1,p_2,p_3\equiv 2\pmod 3$:

\numline{p_1}{6719198563363963000306859976745602011477}
\numline{a_1}{3490701603819271511450476386979369043828}
\numline{p_2}{5014753468477740779475587342292599756573}
\numline{a_2}{2189794512725854244347716309834217468129}
\numline{p_3}{4350736286376564530047381650448847768399}
\numline{a_3}{525176878515189528938215354529508172272}

each individually verified to satisfy $36a_i^3\equiv 19\pmod{p_i}$. The
combined modulus is the 120-digit integer $M=p_1p_2p_3$:

\numline{M}{146598599970416865193503095382122376155797697327993857757772703444309911094348084964151091360030057229085267028040668079}

with CRT representative (119 digits):

\numline{A}{78645806109075982274155955035418384113009237358084530484254201316396650239497999707497613326168293812811504563595780749}

By Theorem~\ref{thm:crt}, for \emph{every} integer $n\equiv A\pmod M$,
\[
  x = M \quad\text{exactly divides}\quad 36n^3-19 .
\]
At $n=A$, $36A^3-19$ is a 359-digit integer, and $(36A^3-19)/M=Q$ is the
following 240-digit integer:

\numline{Q}{119453638517902053794963958283357629777707983376380601474787235288607383318513130562995313813671450899537612665803296173072136255334374568105736108037208680482571587733372866341630091424872047846333343491478053473403358730772112480717305455}

verified exactly via $36A^3-19=M\cdot Q$ in arbitrary-precision arithmetic.

\section{Discussion and honest audit}

\begin{itemize}
\item \textbf{Proved, general, and scalable.} Lemma~\ref{lem:cuberoot} and
Theorem~\ref{thm:crt} are elementary and fully rigorous for \emph{any} choice
of distinct primes $p_i\equiv 2\pmod 3$ (about half of all primes, by
Dirichlet). Nothing here is specific to 30-digit or 40-digit numbers --- the
same construction works verbatim for 100-digit or 1000-digit primes; these
sizes were simply the scales requested. Using $k$ primes of $d$ digits each
yields a modulus $M$ (and hence divisor $x$) of size $\approx dk$ digits.
\item \textbf{What ``largest'' means here.} Since $x$ is forced to be (up to
sign/unit) a product of the $p_i$ for the divisibility to hold across an
entire residue class of $n$ (see \S1), the meaningful notion of ``largest'' is
``most/largest primes multiplied together.'' The $k=3$, 40-digit instance in
\S7 ($x$ a 120-digit number, $n$ a 119-digit number) is the largest computed
here; $k=4,5,\dots$ or larger primes would push this arbitrarily further at
the cost of more/larger primes, with no change in the method.
\item \textbf{If instead a single fixed integer $n$ is wanted} (not a whole
congruence class) with $x\mid 36n^3-19$ for some unrelated, independently-chosen
large $x$, that reduces to integer factorization of $36n^3-19$ for a specific
$n$, which is a fundamentally different (and much harder, non-constructive)
problem --- one has to search and factor rather than construct. The
CRT/cube-root method above sidesteps that difficulty entirely by
\emph{choosing} the primes first and reverse-engineering $n$, which is why it
scales cleanly to 30-digit (or larger) moduli.
\item \textbf{Verification.} Every numerical claim in this document was
checked with exact Python/SymPy arbitrary-precision arithmetic: primality of
each $p_i$, the congruence $p_i\equiv 2\pmod 3$, the root identity
$36a_i^3\equiv 19\pmod{p_i}$, the CRT combination, and the final exact
divisions $36A^3-19 = M\cdot Q$, for the $k=2$ and $k=3$ instances at both the
30-digit and 40-digit scales (four instances total). The closed-form
cube-root formula was additionally cross-checked against brute force for all
primes $p<200$ with $p\equiv2\pmod3$.
\end{itemize}

\end{document}
